\section{A diagonal-completion certificate beyond edgewise payment}

The disjointness of the positive and negative parts permits a second use of
the otherwise invisible diagonal of the cross kernel.  This yields a global
semidefinite certificate even when the two energy kernels are not
$M$-matrices.

\begin{lemma}[Clipped-cross diagonal completion]
\label{lem:psi-diagonal-completion}
Use the zero-constant extensions and symmetric matrices
$K^{\mathrm A},K^{\mathrm B},K^{\mathrm C}$ from
\cref{thm:psi-edgewise-young}.  Let $\widehat K^{\mathrm C}$ be any symmetric
matrix whose off-diagonal entries obey
\begin{equation*}
 \widehat K^{\mathrm C}_{ij}
 \geq\max\{K^{\mathrm C}_{ij},0\},\qquad i\neq j.
\end{equation*}
Its diagonal is arbitrary.  If
\begin{equation}
 \begin{pmatrix}
  K^{\mathrm A}&-\widehat K^{\mathrm C}\\
  -\widehat K^{\mathrm C}&K^{\mathrm B}
 \end{pmatrix}\succeq0,
 \label{eq:psi-diagonal-completion}
\end{equation}
then $\Psi(y,h)\leq0$ for every admissible $y,h$.
\end{lemma}

\begin{proof}
Put $f=y_+$ and $g=y_-$.  They are coordinatewise nonnegative and
$f_ig_i=0$, so the diagonal of the cross matrix does not contribute.  Hence
\begin{equation*}
 f^TK^{\mathrm C}g
 \leq f^T\widehat K^{\mathrm C}g.
\end{equation*}
Applying \cref{eq:psi-diagonal-completion} to $(f,g)$ gives
\begin{equation*}
 2f^TK^{\mathrm C}g
 \leq f^TK^{\mathrm A}f+g^TK^{\mathrm B}g.
\end{equation*}
The eliminated-$h$ identity \cref{eq:psi-eliminated-form} now proves the
claim.
\end{proof}

\begin{theorem}[Twelve- and thirteen-cycle blow-ups]
\label{thm:psi-c12-c13-diagonal-completion}
For $n\in\{12,13\}$, every integer $a\geq1$, and $q=1/20$, the balanced
independent-set blow-up $C_n[\overline K_a]$ satisfies
\begin{equation}
 \sup_{y\in\mathbb R^{na},\ h\in\one^{\perp_D}}\Psi(y,h)=0.
 \label{eq:psi-c12-c13-diagonal-completion-sup}
\end{equation}
On both families condition \textup{(CL)} fails, and neither family is
covered by the edgewise-Young theorem because an off-diagonal entry of each
of $K^{\mathrm A}$ and $K^{\mathrm B}$ is positive.
\end{theorem}

\begin{proof}
After dividing all three $K$-matrices by the common regular degree, exact
rational inversion on each base cycle shows that the only positive
off-diagonal entries of the cross operator $\mathcal C_q$ are the two
adjacent entries.  Write their common value as $c_n>0$ and define the base
operator completion by
\begin{equation*}
 \widehat K^{\mathrm C}_{i,i\pm1}=c_n,
 \qquad
 \widehat K^{\mathrm C}_{ii}=-2c_n,
\end{equation*}
with every other entry zero.  Thus the completion has zero row sums.

The registered exact audit forms the $2n$-by-$2n$ block matrix in
\cref{eq:psi-diagonal-completion}.  Its two constant directions vanish.  In
the rational basis $e_i-e_{n-1}$ in each block, exact $LDL^T$ elimination
has respectively $22$ and $24$ strictly positive pivots.  The smallest
pivots are
\begin{equation*}
 \frac{4298427}{27779000}\quad(n=12),
 \qquad
 \frac{2069613}{14158400}\quad(n=13).
\end{equation*}
This proves the base completion without a floating-point spectral
calculation.

For a balanced blow-up, the operator-level completion is equivalently
\begin{equation*}
 \widehat{\mathcal C}_{\rm full}
 =\widehat{\mathcal C}_{\rm base}\otimes\frac{J_a}{a}
  -2c_n I_n\otimes\left(I_a-\frac{J_a}{a}\right).
\end{equation*}
Thus it equals $c_n/a$ between adjacent parts, $-2c_n$ on every diagonal,
and zero on all other off-diagonal pairs.  Multiplication by the common
degree $2a$ lifts this operator completion to the three $K$-matrices without
changing positivity.  Hence
part-constant vectors inherit exactly the base block, and every actual
off-diagonal cross entry is majorized.  On
each within-part zero-sum space, the three scalar kernel eigenvalues are
obtained by substituting $m=(2/3)m_0$.  The clipped cross completion acts as
$-2c_nI$, and the exact audit verifies
\begin{equation*}
 a_{\rm in}>0,\qquad b_{\rm in}>0,\qquad
 a_{\rm in}b_{\rm in}-(2c_n)^2>0.
\end{equation*}
The actual within-part off-diagonal entry of $K^{\mathrm C}$ is negative, so
no additional clipped edge is missing.  The part-constant and within-part
decomposition therefore proves \cref{eq:psi-diagonal-completion} on every
blow-up size $a$.

Finally, using $3<\pi<22/7$ and
$\sin x\geq x-x^3/6$, the checker proves for both $n=12,13$ that
\begin{equation*}
 2\sin^2(\pi/n)>\frac1{10}=2q.
\end{equation*}
The added within-part modes have normalized-Laplacian eigenvalue one, so the
high-gap hypothesis holds.  Apply
\cref{lem:psi-diagonal-completion,lem:psi-eliminate-h}; equality is attained
at $y=h=0$.
\end{proof}

\section{Current boundary}

Four routes are now analytic: positive association closes every complete
multipartite graph; the clipped-Laplacian sign route closes all balanced
blow-ups of $C_6$ despite strict \textup{(HK)} failure; edgewise Young
payment closes $C_{9},C_{10},C_{11}$ blow-ups at $q=1/20$ (and $C_{10}$ on
an interval); and diagonal completion closes the next two cycle sizes even
though the energy kernels cease to be $M$-matrices.  The next falsifiable
target is a graph-uniform diagonal completion, a quotient-varying
certificate, or an exact admissible positive witness.  A one-step inequality
still does not by itself charge graph discovery, clipping, shifted inverse
work, proper-face changes, finite precision, or final PPR output.
